\subsection{Composición fotónica del nivel \texorpdfstring{\(A_2\)}{A2} de Bendersky--Glaisher}
\label{v:subsec:bendersky-fotonica}

La densidad de número de la proposición~\ref{v:prop:densidades-termicas}
admite una segunda lectura exacta que conserva su procedencia espectral.
Sea \(\mathcal Z_3\) el funcional escalar normalizado producido por el
lector espectral trítico y reconocido posteriormente como la función zeta
en la realización arquimediana. Sean
\[
 H_k=\sum_{j=1}^{k}\frac1j,\qquad H_0=0,
\]
y \(\mathsf B_j\) los números de Bernoulli con el convenio
\(\mathsf B_1=-1/2\). Para evitar toda colisión con la notación de redes,
definimos los niveles de Bendersky--Glaisher mediante
\begin{equation}
 A_k^{\mathrm{BG}}
 =\exp\!\left(
   \frac{H_k\mathsf B_{k+1}}{k+1}-\mathcal Z_3'(-k)
 \right),\qquad k\in\mathbb N_0.
 \label{v:eq:bendersky-bg}
\end{equation}
El término de Bernoulli fija la normalización del producto regularizado;
no introduce una constante física. La coordenada
\(\pi_{\mathrm{HMT}}\) que interviene en la ecuación funcional ya fue
expresada por el estado enriquecido APP--TRIT--TPK antes de esta lectura.

\begin{proposition}[Composición fotónica del nivel dos]
\label{v:prop:bendersky-fotonica}
En la normalización escalar \(\mathcal Z_3=\zeta\),
\begin{equation}
 \log A_2^{\mathrm{BG}}
 =-\mathcal Z_3'(-2)
 =\frac{\mathcal Z_3(3)}{4\pi_{\mathrm{HMT}}^2}.
 \label{v:eq:bendersky-nivel-dos}
\end{equation}
Por consiguiente, en la coordenatización térmica del calendario,
\begin{align}
 n_\gamma\ell_P^3
  &=\frac{8\log A_2^{\mathrm{BG}}}{(\log3)^3}
   =-\frac{8\mathcal Z_3'(-2)}{(\log3)^3},
   \label{v:eq:numero-fotones-bendersky}\\
 \frac{u_\gamma}{n_\gamma k_B\Theta_{108}}
  &=\frac{\pi_{\mathrm{HMT}}^2}{120\log A_2^{\mathrm{BG}}},
   \label{v:eq:energia-media-bendersky}\\
 \frac{s_\gamma}{n_\gamma k_B}
  &=\frac{\pi_{\mathrm{HMT}}^2}{90\log A_2^{\mathrm{BG}}}.
   \label{v:eq:entropia-media-bendersky}
\end{align}
\end{proposition}

\begin{proof}
La forma completada satisface
\[
 \pi_{\mathrm{HMT}}^{-s/2}\Gamma(s/2)\mathcal Z_3(s)
 =\pi_{\mathrm{HMT}}^{-(1-s)/2}
  \Gamma((1-s)/2)\mathcal Z_3(1-s).
\]
En \(s=-2\), el cero simple de \(\mathcal Z_3\) cancela el polo de
\(\Gamma(s/2)\). Comparar los términos constantes, usando
\(\Gamma(3/2)=\sqrt{\pi_{\mathrm{HMT}}}/2\) tras el reconocimiento
arquimediano, da
\[
 \mathcal Z_3'(-2)=-\frac{\mathcal Z_3(3)}{4\pi_{\mathrm{HMT}}^2}.
\]
Como \(\mathsf B_3=0\), la definición~\eqref{v:eq:bendersky-bg}
produce la primera igualdad de~\eqref{v:eq:bendersky-nivel-dos}.
Sustituir
\(\mathcal Z_3(3)=4\pi_{\mathrm{HMT}}^2\log A_2^{\mathrm{BG}}\)
en~\eqref{v:eq:densidades-calendario} demuestra
\eqref{v:eq:numero-fotones-bendersky}. Finalmente,
\(k_B\Theta_{108}=E_P/\log3\); dividir las densidades de energía y
entropía por la densidad de número da
\eqref{v:eq:energia-media-bendersky} y
\eqref{v:eq:entropia-media-bendersky}.
\end{proof}

La evaluación \(A_2^{\mathrm{BG}}=1.0309167521973921\ldots\) sirve como
test numérico de las tres identidades, no como dato generador. El desarrollo
espectral Bendersky--Glaisher contiene la jerarquía completa de derivadas
regularizadas y su correspondencia con los valores espectrales impares; la presente sección
establece la composición recíproca que utiliza específicamente la densidad
adimensional de número fotónico. No se redefine \(k_B\), no se identifica una densidad con un
operador regularizado y no se altera ninguna de las densidades con
\(\zeta(3)\) ya demostradas.
